\documentclass[11pt]{article}
\usepackage{amsmath,amsthm,amssymb}
\usepackage{mathpartir}
\usepackage{stmaryrd}
\usepackage[margin=1in]{geometry}

\newtheorem{theorem}{Theorem}
\newtheorem{lemma}[theorem]{Lemma}
\newtheorem{definition}{Definition}
\newtheorem{corollary}[theorem]{Corollary}

\title{Causal Programming: do-Calculus as a Type System Construct}
\author{Sounio Language Research Team}
\date{Draft --- \today}

\begin{document}
\maketitle

\begin{abstract}
We present a type system extension that embeds Pearl's do-calculus into the programming language's type system, enabling \emph{compile-time verification} of causal identifiability. Our system encodes causal graphs as types, the do-operator as a type constructor, and d-separation as a type-level predicate verified by SMT solvers. We prove that well-typed causal programs cannot produce spurious correlations: if the type checker accepts a do-query, the causal effect is identifiable from observational data. This is the first programming language with compile-time causal reasoning guarantees.
\end{abstract}

\section{Introduction}

Causal inference is fundamental to science, medicine, and policy. Yet mainstream programming languages treat causal reasoning as untyped computation:
\begin{itemize}
\item \textbf{Confounding bugs}: Programmers forget to condition on confounders
\item \textbf{Non-identifiable queries}: Computing ``causal effects'' that cannot be identified from data
\item \textbf{Spurious correlations}: Confusing association with causation
\end{itemize}

\noindent We introduce \textbf{causal types} that solve these problems by:
\begin{enumerate}
\item Encoding causal DAGs in the type system
\item Implementing the do-operator as a type constructor
\item Verifying identifiability at compile time via SMT
\item Integrating with epistemic types for uncertainty quantification
\end{enumerate}

\subsection{Pearl's Causal Hierarchy}

Our type system encodes all three levels of Pearl's ladder of causation:

\begin{enumerate}
\item \textbf{Level 1 --- Association}: $P(Y \mid X)$ --- Standard types
\item \textbf{Level 2 --- Intervention}: $P(Y \mid \text{do}(X))$ --- Causal types
\item \textbf{Level 3 --- Counterfactual}: $P(Y_x \mid X=x', Y=y)$ --- Structural types
\end{enumerate}

\section{Syntax}

\subsection{Causal Graph Types}
\[
\begin{array}{rcl}
\mathcal{G} & ::= & (\mathcal{V}, \mathcal{E}, \mathcal{B}) \\
\mathcal{V} & ::= & \{V_1, \ldots, V_n\} \quad \text{(variable nodes)} \\
\mathcal{E} & \subseteq & \mathcal{V} \times \mathcal{V} \quad \text{(directed edges)} \\
\mathcal{B} & \subseteq & \binom{\mathcal{V}}{2} \quad \text{(bidirected edges, latent confounders)}
\end{array}
\]

\subsection{Causal Types}
\[
\begin{array}{rcl}
\tau & ::= & \ldots \quad \text{(standard types)} \\
    & \mid & \texttt{CausalEffect}\langle T, \mathcal{G}, X, Y \rangle \\
    & \mid & \texttt{CausalKnowledge}\langle T, \epsilon, \mathcal{G}, \text{do}(X) \rangle \\
    & \mid & \texttt{Graph}\langle \mathcal{G} \rangle
\end{array}
\]

\subsection{Causal Terms}
\[
\begin{array}{rcl}
e & ::= & \ldots \quad \text{(standard terms)} \\
  & \mid & \texttt{graph}\{ V_1 \to V_2, \ldots \} \quad \text{(graph literal)} \\
  & \mid & \texttt{do}(\mathcal{G}, X = x) \quad \text{(intervention)} \\
  & \mid & \texttt{identify}(\mathcal{G}, X, Y) \quad \text{(identification query)} \\
  & \mid & \texttt{adjust}(\mathcal{G}, X, Y, Z) \quad \text{(adjustment formula)}
\end{array}
\]

\section{Type-Level Causal Predicates}

\subsection{d-Separation}

\begin{definition}[d-Separation]
A path $\pi$ between $X$ and $Y$ in $\mathcal{G}$ is \emph{blocked} by a set $Z$ if there exists a node $W$ on $\pi$ such that either:
\begin{enumerate}
\item $W$ is a non-collider ($\to W \to$ or $\gets W \to$) and $W \in Z$, or
\item $W$ is a collider ($\to W \gets$) and neither $W$ nor its descendants are in $Z$.
\end{enumerate}
$X$ and $Y$ are \emph{d-separated} given $Z$ (written $X \perp_\mathcal{G} Y \mid Z$) if every path between them is blocked by $Z$.
\end{definition}

We encode d-separation as a type-level predicate:
\[
\texttt{d\_separated}(\mathcal{G}, X, Y, Z) : \texttt{Prop}
\]
verified by the SMT solver at compile time.

\subsection{Backdoor Criterion}

\begin{definition}[Backdoor Criterion]
A set $Z$ satisfies the backdoor criterion relative to $(X, Y)$ in $\mathcal{G}$ if:
\begin{enumerate}
\item No node in $Z$ is a descendant of $X$
\item $Z$ blocks every path between $X$ and $Y$ that contains an arrow into $X$
\end{enumerate}
\end{definition}

\[
\texttt{backdoor}(\mathcal{G}, X, Y, Z) \triangleq
Z \cap \text{desc}(X) = \emptyset \;\wedge\;
X \perp_{\mathcal{G}_{\overline{X}}} Y \mid Z
\]

\subsection{Frontdoor Criterion}

\begin{definition}[Frontdoor Criterion]
A set $M$ satisfies the frontdoor criterion relative to $(X, Y)$ in $\mathcal{G}$ if:
\begin{enumerate}
\item $M$ intercepts all directed paths from $X$ to $Y$
\item There is no unblocked back-door path from $X$ to $M$
\item All back-door paths from $M$ to $Y$ are blocked by $X$
\end{enumerate}
\end{definition}

\subsection{Instrumental Variable Criterion}

\begin{definition}[Instrumental Variable]
A variable $Z$ is an \emph{instrument} for the causal effect of $X$ on $Y$ in $\mathcal{G}$ if:
\begin{enumerate}
\item $Z$ affects $X$: there is a directed path from $Z$ to $X$
\item $Z$ does not directly affect $Y$: $Z$ and $Y$ are d-separated in $\mathcal{G}_{\overline{X}}$
\item $Z$ and $Y$ share no common causes: no confounding path $Z \gets \cdots \to Y$
\end{enumerate}
\end{definition}

\[
\texttt{instrument}(\mathcal{G}, Z, X, Y) \triangleq
Z \to^* X \;\wedge\;
Z \perp_{\mathcal{G}_{\overline{X}}} Y \;\wedge\;
\lnot\exists U.\; Z \gets U \to^* Y
\]

\subsection{Identifiability}
\begin{definition}[Identifiability]
The causal effect $P(Y \mid \text{do}(X))$ is \emph{identifiable} from $\mathcal{G}$ if it can be uniquely computed from the observational distribution $P(\mathcal{V})$ for any positive distribution compatible with $\mathcal{G}$.
\end{definition}

\[
\texttt{identifiable}(\mathcal{G}, X, Y) \triangleq
\exists Z.\; \texttt{backdoor}(\mathcal{G}, X, Y, Z) \;\lor\;
\exists M.\; \texttt{frontdoor}(\mathcal{G}, X, Y, M) \;\lor\;
\exists I.\; \texttt{instrument}(\mathcal{G}, I, X, Y) \;\lor\;
\texttt{id\_algorithm}(\mathcal{G}, X, Y) = \top
\]

\section{Typing Rules}

\subsection{Graph Construction}
\[
\inferrule*[right=T-Graph]
  {\mathcal{G} \text{ is a valid DAG} \\
   \forall (V_i, V_j) \in \mathcal{E}.\; V_i \neq V_j \\
   \text{no directed cycles in } \mathcal{G}}
  {\Gamma \vdash \texttt{graph}\{\ldots\} : \texttt{Graph}\langle \mathcal{G} \rangle}
\]

\subsection{do-Operator}
\[
\inferrule*[right=T-Do]
  {\Gamma \vdash g : \texttt{Graph}\langle \mathcal{G} \rangle \\
   X \in \mathcal{V}(\mathcal{G}) \quad Y \in \mathcal{V}(\mathcal{G}) \\
   \texttt{SMT-PROVE}(\texttt{identifiable}(\mathcal{G}, X, Y))}
  {\Gamma \vdash \texttt{do}(g, X = x) : \texttt{CausalEffect}\langle T_Y, \mathcal{G}, X, Y \rangle}
\]

\noindent \textbf{Key property}: The typing rule \emph{requires} that identifiability is proven before the do-operator can be used. If the SMT solver cannot prove identifiability, the program is rejected at compile time.

\subsection{Backdoor Adjustment}
\[
\inferrule*[right=T-Backdoor]
  {\Gamma \vdash g : \texttt{Graph}\langle \mathcal{G} \rangle \\
   \texttt{SMT-PROVE}(\texttt{backdoor}(\mathcal{G}, X, Y, Z)) \\
   \Gamma \vdash \texttt{data} : \texttt{Data}\langle \{X, Y\} \cup Z \rangle}
  {\Gamma \vdash \texttt{adjust}(g, X, Y, Z) : \texttt{CausalEffect}\langle T_Y, \mathcal{G}, X, Y \rangle}
\]

\subsection{Instrumental Variable Estimation}
\[
\inferrule*[right=T-IV]
  {\Gamma \vdash g : \texttt{Graph}\langle \mathcal{G} \rangle \\
   \texttt{SMT-PROVE}(\texttt{instrument}(\mathcal{G}, Z, X, Y)) \\
   \Gamma \vdash \texttt{data} : \texttt{Data}\langle \{X, Y, Z\} \rangle}
  {\Gamma \vdash \texttt{iv\_estimate}(g, Z, X, Y) : \texttt{CausalEffect}\langle T_Y, \mathcal{G}, X, Y \rangle}
\]

\subsection{Integration with Epistemic Types}
\[
\inferrule*[right=T-CausalEpistemic]
  {\Gamma \vdash e : \texttt{CausalEffect}\langle T, \mathcal{G}, X, Y \rangle \\
   \epsilon_{\text{stat}} = \text{statistical\_uncertainty}(e) \\
   \epsilon_{\text{causal}} = \text{positivity\_uncertainty}(\mathcal{G}, X) \\
   \epsilon = \sqrt{\epsilon_{\text{stat}}^2 + \epsilon_{\text{causal}}^2}}
  {\Gamma \vdash \texttt{estimate}(e) : \texttt{Knowledge}\langle T, \epsilon, \texttt{Causal}, t \rangle}
\]

\noindent The causal uncertainty $\epsilon_{\text{causal}}$ captures potential violations of positivity, while $\epsilon_{\text{stat}}$ captures finite-sample uncertainty. Both are tracked in the epistemic type.

\section{Operational Semantics}

\subsection{Causal Values}
\[
v_c ::= \langle v, \mathcal{G}, X, Y, Z_{\text{adj}} \rangle
\]
where $v$ is the estimated effect, $\mathcal{G}$ is the causal graph, and $Z_{\text{adj}}$ is the adjustment set used.

\subsection{Reduction Rules}

\[
\inferrule*[right=E-Do-Backdoor]
  {\texttt{backdoor}(\mathcal{G}, X, Y, Z) \\
   Z_{\text{adj}} = \text{find\_adjustment}(\mathcal{G}, X, Y)}
  {\texttt{do}(\mathcal{G}, X = x) \longrightarrow
   \sum_{z} P(Y \mid X = x, Z = z) \cdot P(Z = z)}
\]

\[
\inferrule*[right=E-Do-Frontdoor]
  {\texttt{frontdoor}(\mathcal{G}, X, Y, M) \\
   M = \text{find\_mediators}(\mathcal{G}, X, Y)}
  {\texttt{do}(\mathcal{G}, X = x) \longrightarrow
   \sum_{m} P(M = m \mid X = x) \sum_{x'} P(Y \mid M = m, X = x') P(X = x')}
\]

\[
\inferrule*[right=E-Do-IV]
  {\texttt{instrument}(\mathcal{G}, Z, X, Y) \\
   \hat{\beta}_{\text{IV}} = \text{Cov}(Y, Z) / \text{Cov}(X, Z)}
  {\texttt{iv\_estimate}(\mathcal{G}, Z, X, Y) \longrightarrow
   \hat{\beta}_{\text{IV}}}
\]

\section{Metatheory}

\subsection{Causal Soundness}
\begin{theorem}[Causal Soundness]
If $\Gamma \vdash e : \texttt{CausalEffect}\langle T, \mathcal{G}, X, Y \rangle$, then the computed effect equals the interventional distribution $P(Y \mid \text{do}(X))$, not merely the observational conditional $P(Y \mid X)$.
\end{theorem}

\begin{proof}[Proof Sketch]
By the typing rule T-Do, the program type-checks only if $\texttt{identifiable}(\mathcal{G}, X, Y)$ is proven by the SMT solver. By the completeness of the ID algorithm \cite{tian2002general}, identifiability implies that $P(Y \mid \text{do}(X))$ can be uniquely computed from observational data. The operational semantics (E-Do-Backdoor, E-Do-Frontdoor) implement exactly the adjustment formulas guaranteed to be correct by the backdoor and frontdoor criteria. Therefore, the computed value equals the true interventional distribution.
\end{proof}

\subsection{Non-Identifiability Rejection}
\begin{theorem}[Non-Identifiable Rejection]
If the causal effect $P(Y \mid \text{do}(X))$ is not identifiable from $\mathcal{G}$, then no well-typed program can compute it.
\end{theorem}

\begin{proof}
The typing rule T-Do requires $\texttt{SMT-PROVE}(\texttt{identifiable}(\mathcal{G}, X, Y))$. If the effect is non-identifiable, no valid proof exists, so the SMT solver returns Unknown or Unsat. The type checker then rejects the program with a compile-time error, possibly providing a counterexample (a ``hedge'' demonstrating non-identifiability).
\end{proof}

\subsection{Progress and Preservation}

\begin{theorem}[Progress]
If $\emptyset \vdash e : \texttt{CausalEffect}\langle T, \mathcal{G}, X, Y \rangle$, then either $e$ is a causal value or there exists $e'$ such that $e \longrightarrow e'$.
\end{theorem}

\begin{theorem}[Preservation]
If $\Gamma \vdash e : \texttt{CausalEffect}\langle T, \mathcal{G}, X, Y \rangle$ and $e \longrightarrow e'$, then $\Gamma \vdash e' : \texttt{CausalEffect}\langle T', \mathcal{G}, X, Y \rangle$ where $T'$ is a refinement of $T$.
\end{theorem}

\section{Counterfactual Types (Level 3)}

Counterfactual reasoning asks ``What would $Y$ have been, had $X$ been $x$, given that we observed $X = x'$?'' We encode this as:

\[
\texttt{Counterfactual}\langle T, \mathcal{G}, X{=}x, \text{evidence} \rangle
\]

\subsection{Typing Rule}
\[
\inferrule*[right=T-CF]
  {\Gamma \vdash g : \texttt{Graph}\langle \mathcal{G} \rangle \\
   \mathcal{G} \text{ is a recursive structural model} \\
   \Gamma \vdash \text{evidence} : \texttt{Data}\langle \{X = x', Y = y\} \rangle}
  {\Gamma \vdash \texttt{cf}(g, X{=}x, \text{evidence}) : \texttt{Counterfactual}\langle T_Y, \mathcal{G}, X{=}x, \text{evidence} \rangle}
\]

\noindent Counterfactual evaluation requires the full structural causal model, not just the graph. The typing rule ensures the SCM is available.

\subsection{Three-Step Algorithm}
The operational semantics for counterfactuals follows Pearl's three-step procedure:
\begin{enumerate}
\item \textbf{Abduction:} Compute $P(U \mid \text{evidence})$ — update exogenous variables
\item \textbf{Action:} Modify the structural equations to enforce $X = x$
\item \textbf{Prediction:} Compute $Y$ under the modified model with updated exogenous variables
\end{enumerate}

\[
\inferrule*[right=E-CF]
  {U^* = \text{abduct}(\mathcal{M}, \text{evidence}) \\
   \mathcal{M}_{X=x} = \text{intervene}(\mathcal{M}, X{=}x) \\
   v = \text{predict}(\mathcal{M}_{X=x}, U^*)}
  {\texttt{cf}(g, X{=}x, \text{evidence}) \longrightarrow v}
\]

\section{Integration with Epistemic Types}

\subsection{CausalKnowledge Type}
The full integration of causal and epistemic types yields:
\[
\texttt{CausalKnowledge}\langle T, \epsilon, \Phi, t, \mathcal{G}, \text{do}(X) \rangle
\]
which carries:
\begin{itemize}
\item $T$: the value type (e.g., Average Treatment Effect)
\item $\epsilon$: combined uncertainty ($\epsilon_{\text{stat}} \oplus \epsilon_{\text{causal}}$)
\item $\Phi$: provenance (including which adjustment formula was used)
\item $t$: temporal validity
\item $\mathcal{G}$: the causal graph used
\item $\text{do}(X)$: the intervention performed
\end{itemize}

\subsection{Uncertainty Decomposition}
For a causal estimate, we can decompose uncertainty into components:
\[
\epsilon_{\text{total}}^2 = \epsilon_{\text{stat}}^2 + \epsilon_{\text{causal}}^2 + \epsilon_{\text{model}}^2
\]
where:
\begin{itemize}
\item $\epsilon_{\text{stat}}$: finite-sample uncertainty (decreases with $n$)
\item $\epsilon_{\text{causal}}$: positivity and overlap uncertainty
\item $\epsilon_{\text{model}}$: structural model misspecification
\end{itemize}

\section{Case Studies}

\subsection{PBPK Pharmacokinetic Model}
\begin{verbatim}
let g = CausalGraph::new([
    edge("Dose", "Plasma"),
    edge("Plasma", "Effect"),
    edge("Genotype", "Metabolism"),
    edge("Metabolism", "Plasma")
])

// Type checks: identifiable via backdoor {Genotype, Metabolism}
let effect: CausalKnowledge<Concentration> =
    g.do(Dose = 100.0 mg)
\end{verbatim}

\subsection{fMRI Causal Connectivity}
\begin{verbatim}
let brain = CausalGraph::from_atlas(AAL_116)

// Compile-time error: V1 -> MT not identifiable
// (unobserved confounder between V1 and MT)
let bad: CausalEffect<f64> =
    brain.do(V1 = active)  // TYPE ERROR!
\end{verbatim}

\section{Related Work}

\paragraph{DoWhy/EconML} Python libraries for causal inference, but no compile-time verification.

\paragraph{Ananke} Causal inference in R with graphical criteria, but runtime checks only.

\paragraph{Our Contribution} First \emph{type system} with:
\begin{itemize}
\item Compile-time identifiability verification via SMT
\item do-operator as a type constructor (not just a function)
\item Causal-epistemic integration (uncertainty decomposition)
\item Provable rejection of non-identifiable queries
\end{itemize}

\section{Conclusion}

Causal types bring Pearl's do-calculus into the type system, preventing confounding bugs at compile time. Combined with epistemic types, this enables fully tracked causal-epistemic reasoning where every causal estimate carries its uncertainty provenance and identifiability proof.

Future work includes:
\begin{itemize}
\item Counterfactual types (Level 3 of the causal hierarchy)
\item Sensitivity analysis types (bounds under violation of assumptions)
\item Integration with randomized experiment types
\end{itemize}

\begin{thebibliography}{9}

\bibitem{pearl2009causality}
J.~Pearl.
\newblock {\em Causality: Models, Reasoning, and Inference}.
\newblock Cambridge University Press, 2nd edition, 2009.

\bibitem{tian2002general}
J.~Tian and J.~Pearl.
\newblock On the identification of causal effects.
\newblock In {\em UAI}, 2002.

\bibitem{angrist1996identification}
J.~D. Angrist, G.~W. Imbens, and D.~B. Rubin.
\newblock Identification of causal effects using instrumental variables.
\newblock {\em Journal of the American Statistical Association}, 91(434):444--455, 1996.

\bibitem{pearl1995causal}
J.~Pearl.
\newblock Causal diagrams for empirical research.
\newblock {\em Biometrika}, 82(4):669--688, 1995.

\bibitem{bareinboim2016causal}
E.~Bareinboim and J.~Pearl.
\newblock Causal inference and the data-fusion problem.
\newblock {\em PNAS}, 113(27):7345--7352, 2016.

\end{thebibliography}

\end{document}
